\section{Related work}\label{sec:related}

\subsection{Classical and dynamic scheduling}
Single-resource sequencing theory provides the canonical dispatching rules and their optimality conditions: SPT for mean flow time \cite{smith1956}, EDF for maximum lateness \cite{jackson1955}, and the Moore--Hodgson algorithm for the number of late jobs \cite{moore1968}. The three-field classification of \citet{graham1979} and the textbook treatment of \citet{pinedo2016} organise the deterministic theory. In real-time computing, EDF is optimal for preemptive scheduling of independent tasks whenever the processor is not overloaded \cite{liu1973}, but its performance can collapse under overload, when a domino effect causes many deadlines to be missed \cite{buttazzo2011}. This failure mode is directly relevant here, because a fatigued worker is in effect an overloaded resource. Surveys of priority rules \cite{panwalkar1977,haupt1989} and of dynamic and reactive scheduling \cite{vieira2003,ouelhadj2009} show that practice relies heavily on such rules when information arrives online. In all of these models, however, processing times are exogenous: they do not depend on the history of the schedule.

\subsection{Human factors in scheduling and workforce models}
Personnel scheduling and rostering \cite{ernst2004,burke2004,vandenbergh2013} assign workers to shifts but typically treat performance within a shift as constant. Research at the interface of operations and human resources argues that this abstraction misses first-order effects \cite{boudreau2003,neumann2010}, and \citet{lodree2009} proposed a taxonomy for integrating human factors into scheduling theory. Quantitative models of learning, forgetting, fatigue and recovery have been developed for dual-resource-constrained systems \cite{wright1936,jaber2010,jaber2013,givi2015}. They use exponential fatigue accumulation during work and exponential recovery during rest, which is the functional form we adopt. Rest-break scheduling has been formulated as an optimisation problem since \citet{bechtold1984}, with ergonomic guidance on work--rest ratios summarised by \citet{konz1998} and rest-allowance models for task assignment developed by \citet{calzavara2019}. Fatigue has also been incorporated into flexible job-shop formulations solved by metaheuristics \cite{tan2021} and, very recently, by graph-based DRL \cite{hu2025}. Reviews of human factors in Industry~4.0 and in order picking \cite{grosse2017,kadir2019,neumann2021,katiraee2021} conclude that worker heterogeneity and well-being are still under-represented in operational models. A common feature of this literature is that the worker state is assumed to be known, either deterministic or computed from the schedule, so the decision problem is fully observable.

\subsection{Deep reinforcement learning for scheduling}
Reinforcement learning has been applied to scheduling since the mid-1990s \cite{zhang1995,aydin2000}, and the combination of Q-learning \cite{watkins1992} with deep function approximation \cite{mnih2015} enabled scheduling policies that operate on rich state representations. Notable examples include cluster resource management \cite{mao2016}, data-processing job scheduling with graph neural networks \cite{mao2019}, semiconductor production scheduling \cite{waschneck2018}, dynamic flexible job shops with job insertions \cite{luo2020,song2023}, and size-agnostic job-shop dispatching \cite{zhang2020l2d,park2021,liu2020}. Parameter-adaptive RL for dynamic job shops \cite{shahrabi2017} and RL for combinatorial optimisation more broadly \cite{bengio2021,mazyavkina2021} complete the picture. Systematic reviews \cite{kayhan2023,panzer2022,khadivi2025} note three recurring weaknesses: comparisons against untuned or weak baselines, limited statistical analysis, and little attention to human resources. Our study addresses all three.

\subsection{Cognitive state, its measurement, and adaptive systems}
Resource theories of attention \cite{kahneman1973,wickens2008} and cognitive load theory \cite{sweller1988} model mental capacity as a limited resource. Vigilance research shows that this resource is depleted by sustained effort \cite{mackworth1948,warm2008} and partially restored by brief breaks \cite{ariga2011}. Accounts of fatigue as an effort-regulation or opportunity-cost signal \cite{hockey2013,kurzban2013} and evidence linking fatigue to disengagement \cite{hopstaken2015} motivate treating fatigue as a state variable that shapes future capacity. Biomathematical fatigue models, from the two-process model of sleep regulation \cite{borbely1982,daan1984,borbely2016} to the three-process model of alertness \cite{akerstedt1997} and applied models such as SAFTE \cite{hursh2004}, combine homeostatic and circadian components. We borrow this decomposition at the time scale of a single workday. Recovery research in organisational psychology confirms the value of micro-breaks \cite{sonnentag2007,fritz2013,kim2017,albulescu2022}. On the measurement side, workload and fatigue can be estimated from EEG, cardiac, ocular and behavioural signals \cite{gevins2003,borghini2014,hogervorst2014,charles2019,shaffer2017}, from vigilance probes such as the PVT \cite{dinges1985,basner2011}, or from questionnaires such as NASA-TLX \cite{hart1988,young2015}. All of these estimates are noisy. Adaptive automation and physiological computing use such estimates to adjust system behaviour \cite{byrne1996,parasuraman2000,kaber2004,fairclough2009,zander2011,feigh2012}, and interruption-management systems decide when to notify a user based on inferred cognitive state \cite{horvitz2003,iqbal2008,czerwinski2004,mark2008}. Reinforcement learning has been used to personalise just-in-time health interventions \cite{klasnja2015,nahumshani2018,liao2020,yu2021}, and in human--robot collaboration to adapt task allocation to human capability and fatigue \cite{gombolay2015,peternel2018,ajoudani2018,lamon2019,zhang2022hrc}. These works show that cognitive-state estimates can drive decisions, but they rarely address the sequential scheduling of a task queue under deadlines.

\subsection{Partial observability and multiple objectives in RL}
When the state is not directly observed, the optimal policy is a function of the belief state \cite{astrom1965,smallwood1973,kaelbling1998}. Memoryless policies can be arbitrarily poor \cite{singh1994}. Practical deep RL approximates the belief either with recurrent networks \cite{hausknecht2015,hochreiter1997,igl2018} or with finite observation histories, as in the frame stacking of \citet{mnih2015}. \citet{ni2022} show that well-tuned memory-based model-free agents are strong POMDP baselines. Stabilisation techniques for DQN, namely experience replay \cite{lin1992}, target networks \cite{mnih2015}, Double Q-learning \cite{vanhasselt2016}, dueling heads \cite{wang2016dueling} and their combinations \cite{hessel2018}, are well established, as is the theory of TD learning with function approximation \cite{tsitsiklis1997}. Scheduling for humans is inherently multi-objective: productivity, deadlines and well-being conflict. Multi-objective RL provides the vocabulary of scalarisation and Pareto fronts \cite{roijers2013,vamplew2011,hayes2022}. Reward shaping must be used with care, because additional reward terms can change the optimal policy \cite{ng1999}. Finally, the deployment of RL in settings that involve people raises safety and reliability concerns \cite{amodei2016,garcia2015,dulacarnold2021}, which we address in Section~\ref{sec:discussion}.

\subsection{Research gap}
Table~\ref{tab:positioning} positions this work against representative streams. To our knowledge, no prior study simultaneously (i) models dynamic attention, fatigue and circadian variation as endogenous state, (ii) treats rest as a scheduling action, (iii) assumes that the cognitive state is only partially observable, (iv) learns an adaptive policy, and (v) benchmarks that policy against tuned, cognition-aware heuristics with multi-seed statistical testing. Our study fills this gap.

\begin{table}[t]
\centering
\caption{Positioning of this study relative to representative research streams (\checkmark: addressed; $\circ$: partially or in a simplified form; --: not addressed).}
\label{tab:positioning}
\small
\setlength{\tabcolsep}{3.5pt}
\begin{tabular}{p{4.6cm}cccccc}
\toprule
Stream (representative works) & \makecell{Dynamic\\human state} & \makecell{Rest as\\action} & \makecell{Partial\\observ.} & \makecell{Adaptive\\policy} & \makecell{Deadlines/\\queue} & \makecell{Tuned cog.\\baselines} \\
\midrule
Dispatching rules \cite{smith1956,jackson1955,liu1973,haupt1989} & -- & -- & -- & -- & \checkmark & -- \\
Fatigue/recovery models in DRC systems \cite{jaber2010,jaber2013,givi2015} & \checkmark & $\circ$ & -- & -- & $\circ$ & -- \\
Rest-break / rest-allowance optimisation \cite{bechtold1984,konz1998,calzavara2019} & $\circ$ & \checkmark & -- & -- & -- & -- \\
DRL for machine scheduling \cite{mao2016,waschneck2018,luo2020,zhang2020l2d,song2023} & -- & -- & -- & \checkmark & \checkmark & -- \\
Fatigue-aware shop scheduling \cite{tan2021,hu2025} & \checkmark & -- & -- & $\circ$ & \checkmark & -- \\
Adaptive automation / interruption mgmt. \cite{byrne1996,horvitz2003,iqbal2008,feigh2012} & \checkmark & $\circ$ & \checkmark & $\circ$ & -- & -- \\
\textbf{This study} & \checkmark & \checkmark & \checkmark & \checkmark & \checkmark & \checkmark \\
\bottomrule
\end{tabular}
\end{table}
